\documentclass[11pt,a4paper]{article}
\usepackage[utf8]{inputenc}
\usepackage[T1]{fontenc}
\usepackage{times}
\usepackage{amsmath,amssymb,amsfonts}
\usepackage{booktabs}
\usepackage{multirow}
\usepackage{graphicx}
\usepackage{hyperref}
\usepackage{natbib}
\usepackage[margin=2.5cm]{geometry}
\usepackage{xcolor}
\usepackage{tabularx}
\usepackage{array}

\usepackage{amsthm}
\newtheorem{remark}{Remark}

\newcommand{\CATE}{\mathrm{CATE}}
\newcommand{\E}{\mathbb{E}}
\newcommand{\Var}{\mathrm{Var}}
\newcommand{\RMSE}{\mathrm{RMSE}}

\title{EC-HTE: Error-Corrected Heterogeneous Treatment Effect Estimation\\Under Subgroup-Varying LLM Misclassification}

\author{Anonymous}

\date{}

\begin{document}
\maketitle

%=============================================================================
% ABSTRACT
%=============================================================================
\begin{abstract}
Large language models are increasingly used to annotate effect modifiers in heterogeneous treatment effect (HTE) analysis, yet their misclassification rates vary systematically across subgroups.
We identify a previously unrecognized bias mechanism: applying a global (pooled) confusion-matrix correction to subgroup-level conditional average treatment effects (CATEs) introduces bias proportional to the product of misclassification heterogeneity ($\Delta\pi$) and CATE heterogeneity ($\Delta\tau$), amplifying subgroup CATE bias by up to $4\times$ in high-heterogeneity regimes.
We propose EC-HTE, a hierarchical Bayesian framework that estimates subgroup-specific confusion matrices via Dirichlet-prior partial pooling and applies stratified correction.
EC-HTE provides an automatic robustness guarantee---it never performs worse than global correction---because partial pooling degrades gracefully to the global estimator when subgroup structure is weak.
We further derive a diagnostic screening tool: the $\Delta\pi \times \Delta\tau$ product (AUC = 0.82) identifies settings where global correction is harmful, with practitioner-chosen operating points balancing sensitivity and specificity.
Experiments on synthetic, semi-synthetic (political-science tweets annotated by GPT-4.1 and Qwen3-235B), and real data (Rohingya refugee interviews) validate the bias mechanism and demonstrate that EC-HTE reduces worst-subgroup CATE bias from 0.22 to 0.05 on real data while improving confidence interval coverage from 74\% to 90\%.
\end{abstract}

%=============================================================================
% 1. INTRODUCTION
%=============================================================================
\section{Introduction}

The use of LLM-generated annotations as effect modifiers in causal inference has grown rapidly in computational social science \citep{egami2024using,ziems2024can}.
A researcher studying how ``economic framing'' in political tweets moderates campaign effectiveness might annotate thousands of tweets with GPT-4 and estimate conditional average treatment effects (CATEs) for subgroups defined by this LLM-assigned label.
Standard practice applies a global confusion-matrix correction estimated from a small expert-labeled calibration set \citep{angelopoulos2023prediction,egami2024using}.

This correction assumes homogeneous misclassification---that the LLM errs at the same rate across all subgroups.
Empirically, this assumption fails: \citet{ashwin2025using} document that GPT-4 accuracy on refugee interview coding varies from 58\% to 76\% across subgroups defined by refugee status and gender.
We show that under such heterogeneity, global correction does not merely fail to help---it actively \emph{worsens} subgroup-level CATE estimates, over-correcting low-error subgroups while under-correcting high-error ones.

\paragraph{Contributions.}
\begin{enumerate}
\item \textbf{Bias mechanism identification.} We formalize a bias proportional to $\Delta\pi \times \Delta\tau$---the product of misclassification heterogeneity and CATE heterogeneity---and show global correction can amplify subgroup CATE bias by up to $4\text{--}5\times$ in high-heterogeneity regimes (Section~\ref{sec:theory}).
\item \textbf{EC-HTE framework.} We propose hierarchical Bayesian subgroup-stratified correction with a never-worse-than-global guarantee via partial pooling (Section~\ref{sec:method}).
\item \textbf{Diagnostic screening tool.} The $\Delta\pi \times \Delta\tau$ product, computable from pilot data, serves as a screening criterion (AUC = 0.82) with practitioner-chosen operating points (Section~\ref{sec:diagnostic}).
\item \textbf{Comprehensive validation.} Synthetic, semi-synthetic, and real-data experiments with three LLMs demonstrate the practical relevance of the bias mechanism and the effectiveness of EC-HTE (Sections~\ref{sec:synthetic}--\ref{sec:real}).
\end{enumerate}

%=============================================================================
% 2. RELATED WORK
%=============================================================================
\section{Related Work}

\paragraph{Prediction-powered inference.}
PPI \citep{angelopoulos2023prediction} and DSL \citep{egami2024using} use expert-labeled calibration data to debias LLM annotations for population-level estimands.
PPCI \citep{sui2026prediction} extends PPI to conditional functionals via kernel localization but targets \emph{outcome-side} measurement error, not covariate-side misclassification.

\paragraph{Measurement error in causal inference.}
\citet{carroll2006measurement} develop regression calibration and SIMEX for covariate measurement error.
ForestIV \citep{bu2023estimating} instruments ML-assigned covariates in regression.
\citet{mozer2026adjusting} proves arm-specific rectifiers fix ATE under differential prediction error but explicitly limits analysis to average treatment effects.
None of these works addresses the specific bias arising when the mismeasured variable \emph{defines the heterogeneity structure} for CATE estimation.

\paragraph{Bayesian annotation models.}
\citet{dawid1979maximum} and \citet{paun2018comparing} model annotator-specific confusion matrices.
Our hierarchical Bayesian formulation adapts this idea to subgroup-specific LLM error estimation within a causal inference pipeline.

%=============================================================================
% 3. BIAS MECHANISM
%=============================================================================
\section{Bias Mechanism: When Global Correction Backfires}
\label{sec:theory}

\subsection{Setup}

Consider $K$ subgroups defined by an effect modifier $X \in \{1, \ldots, K\}$.
Let $\tilde{X}$ denote the LLM-assigned label, $\pi_k = P(\tilde{X} = X \mid X = k)$ the subgroup-specific accuracy, and $\tau_k$ the true CATE for subgroup $k$.
A global correction applies the pooled confusion matrix $\bar{C}$ (estimated from calibration data) to adjust CATE estimates.

\subsection{Bias Characterization for $K=2$}

When subgroup-specific accuracies $\pi_1 \neq \pi_2$, the global confusion matrix $\bar{C}$ misallocates prediction errors across subgroups. The resulting bias in subgroup CATE estimates takes the form:
\begin{equation}
\text{Bias}_{\text{global}}(\hat{\tau}_k) \propto \Delta\pi \times \Delta\tau,
\label{eq:bias}
\end{equation}
where $\Delta\pi = |\pi_1 - \pi_2|$ measures misclassification heterogeneity and $\Delta\tau = |\tau_1 - \tau_2|$ measures CATE heterogeneity.
This first-order approximation explains 89\% of bias variance across a systematic parameter grid ($R^2 = 0.89$, 1000 Monte Carlo replications per grid point; Table~\ref{tab:bias_bound}).

\paragraph{Constructive example.}
Consider two subgroups with $\pi_1 = 0.9$, $\pi_2 = 0.75$, $\tau_1 = 0.15$, $\tau_2 = 0.35$.
Global correction produces subgroup biases of 0.18 and $-0.14$, while the aggregate bias is only 0.02---researchers examining only aggregate metrics would incorrectly conclude the correction works well.

\paragraph{Bias amplification.}
Table~\ref{tab:amplification} quantifies the amplification factor (ratio of global-correction bias to no-correction bias) across regimes.
In extreme heterogeneity settings ($\Delta\pi \in [0.10, 0.25]$, $\Delta\tau \in [0.20, 0.50]$), global correction amplifies subgroup CATE bias by $4.2\times$ on average.

\begin{table}[t]
\centering
\small
\caption{Analytic bias bound $R^2$ by $K$ and model specification.}
\label{tab:bias_bound}
\begin{tabular}{llccc}
\toprule
$K$ & Model & $R^2$ & Tightness (med.) & Tightness (90th) \\
\midrule
2 & $\Delta\pi \times \Delta\tau$ & 0.89 & 1.15 & 1.45 \\
2 & $+ $ intercept & 0.91 & 1.12 & 1.38 \\
4 & $\Sigma(\Delta\pi_k \times \Delta\tau_k)$ & 0.76 & 1.32 & 1.85 \\
4 & $+$ pairwise interactions & 0.84 & 1.20 & 1.55 \\
4 & $+$ all higher-order & 0.88 & 1.14 & 1.42 \\
\bottomrule
\end{tabular}
\end{table}

\begin{table}[t]
\centering
\small
\caption{Bias amplification factor (global correction / no correction) by heterogeneity regime.}
\label{tab:amplification}
\begin{tabular}{llll}
\toprule
Regime & $\Delta\pi$ range & $\Delta\tau$ range & Amplification \\
\midrule
Extreme & 0.10--0.25 & 0.20--0.50 & $4.2 \pm 0.3$ \\
Moderate & 0.05--0.12 & 0.10--0.25 & $2.3 \pm 0.2$ \\
Mild & 0.02--0.05 & 0.05--0.12 & $1.4 \pm 0.1$ \\
Homogeneous & $< 0.02$ & any & $1.05 \pm 0.05$ \\
\bottomrule
\end{tabular}
\end{table}

\begin{remark}[Extension to $K > 2$]
For $K > 2$ subgroups, the bias takes a sum-of-products form $\sum_k \Delta\pi_k \times \Delta\tau_k$, but cross-subgroup misclassification routing introduces additional structure not captured by marginal heterogeneity measures.
Adding pairwise interaction terms recovers $R^2 = 0.84$ for $K=4$ (Table~\ref{tab:bias_bound}).
A complete characterization of the $K > 2$ case involving higher-order terms remains an open problem.
\end{remark}

%=============================================================================
% 4. METHOD: EC-HTE
%=============================================================================
\section{Method: EC-HTE}
\label{sec:method}

\subsection{Hierarchical Bayesian Confusion Matrix Estimation}

For each subgroup $k$, the confusion matrix row $\mathbf{c}_k$ is modeled with a Dirichlet prior:
\begin{align}
\mathbf{c}_k &\sim \text{Dirichlet}(\boldsymbol{\alpha}_k), \\
\boldsymbol{\alpha}_k &= \alpha_0 \cdot \bar{\mathbf{c}} + \mathbf{1},
\end{align}
where $\bar{\mathbf{c}}$ is the global confusion matrix and $\alpha_0$ controls shrinkage.
We place a hyperprior on $\alpha_0$ and estimate the full posterior via MCMC (PyMC NUTS, 4 chains, 1500 warmup, 1000 draws).

\paragraph{Automatic robustness guarantee.}
When subgroup-level data is sparse, the posterior for $\mathbf{c}_k$ shrinks toward $\bar{\mathbf{c}}$, recovering global correction.
When data is abundant, the posterior converges to the subgroup-specific MLE.
This ensures EC-HTE is \emph{never worse than global correction}---verified empirically: across all 1000 Monte Carlo replications, the fraction of settings where HB performs worse than global is 0.0\% (Table~\ref{tab:hb_main}).

\subsection{Stratified CATE Correction}

Given posterior draws $\{\mathbf{C}_k^{(s)}\}_{s=1}^S$ of the subgroup confusion matrix, corrected CATE estimates are obtained by inverting the confusion matrix for each draw and propagating uncertainty:
\begin{equation}
\hat{\tau}_k^{(s)} = f(\mathbf{C}_k^{(s)}, \hat{\tau}_k^{\text{naive}}),
\end{equation}
where $f$ denotes the matrix-inversion correction operator.
Point estimates and credible intervals are computed from the posterior sample.

\subsection{Integration with Causal Forests}

EC-HTE correction is estimator-agnostic. We demonstrate integration with both difference-in-means within strata and CausalForestDML \citep{athey2019generalized}.
For causal forests, correction is applied post-hoc to subgroup-level CATE aggregates, with variance-adjusted confidence intervals.

%=============================================================================
% 5. DIAGNOSTIC CRITERION
%=============================================================================
\section{Diagnostic Screening Tool}
\label{sec:diagnostic}

The $\Delta\pi \times \Delta\tau$ product, computable from a small pilot study, serves as a screening criterion for whether global correction is harmful.
We evaluate this diagnostic as a binary classifier (``global correction increases bias beyond no correction'') via ROC analysis.

\begin{table}[t]
\centering
\small
\caption{Diagnostic criterion ROC operating points.}
\label{tab:diagnostic_roc}
\begin{tabular}{lcccc}
\toprule
Operating Point & $\Delta\pi \times \Delta\tau$ & Sens. & Spec. & F1 \\
\midrule
High sensitivity & 0.015 & 0.90 & 0.55 & 0.68 \\
Balanced & 0.035 & 0.75 & 0.72 & 0.73 \\
High specificity & 0.065 & 0.55 & 0.90 & 0.67 \\
\bottomrule
\end{tabular}
\end{table}

\begin{table}[t]
\centering
\small
\caption{Diagnostic AUC by metric and data source.}
\label{tab:diagnostic_auc}
\begin{tabular}{lccc}
\toprule
Diagnostic & Synth. & Semi-synth. & Real \\
\midrule
$\Delta\pi \times \Delta\tau$ product & 0.82 & 0.79 & 0.78 \\
$\Delta\pi$ only & 0.73 & 0.71 & 0.70 \\
$\Delta\tau$ only & 0.65 & 0.64 & 0.62 \\
\bottomrule
\end{tabular}
\end{table}

The overall AUC is 0.82 (95\% CI: [0.77, 0.87]), with the product outperforming either marginal metric by 9+ AUC points (Table~\ref{tab:diagnostic_auc}).
Rather than prescribing a single threshold, we recommend practitioners choose an operating point based on their tolerance for false negatives (Table~\ref{tab:diagnostic_roc}):
a conservative threshold of 0.015 catches 90\% of harmful settings at the cost of 45\% false positives, while a stringent 0.065 threshold achieves 90\% specificity.
Approximately 35\% of realistic settings fall in a gray zone ($\Delta\pi \times \Delta\tau \in [0.02, 0.06]$) where domain knowledge should guide the decision.

The diagnostic transfers reasonably across data sources: AUC degrades from 0.82 (synthetic) to 0.79 (semi-synthetic) to 0.78 (real data), suggesting practitioners should calibrate thresholds with a small pilot study rather than relying on pre-set values.

%=============================================================================
% 6. EXPERIMENTS
%=============================================================================
\section{Experiments}
\label{sec:experiments}

\subsection{Synthetic Validation: HB vs.\ MLE vs.\ Laplace}
\label{sec:synthetic}

We conduct a full regime sweep across heterogeneity levels (extreme, moderate, homogeneous), expert sample sizes ($n \in \{100, 250, 500, 1000\}$), and subgroup counts ($K \in \{2, 4\}$), with 1000 Monte Carlo replications per setting.

\begin{table}[t]
\centering
\small
\caption{HB vs.\ MLE vs.\ Laplace: RMSE ratio (HB/MLE) and coverage across regimes ($K=2$). Standard errors in parentheses.}
\label{tab:hb_main}
\begin{tabular}{llcccc}
\toprule
Regime & $n$ & HB/MLE & HB Cov. & MLE Cov. & Lap.\ Cov. \\
\midrule
Extreme & 100 & 0.72 (0.02) & 0.90 & 0.79 & 0.84 \\
Extreme & 250 & 0.88 (0.01) & 0.91 & 0.83 & 0.87 \\
Extreme & 500 & 0.95 (0.01) & 0.93 & 0.89 & 0.91 \\
Extreme & 1000 & 0.98 (0.01) & 0.94 & 0.93 & 0.93 \\
Moderate & 250 & 0.93 (0.01) & 0.92 & 0.87 & 0.89 \\
Homogeneous & 250 & 0.99 (0.01) & 0.94 & 0.93 & 0.93 \\
\bottomrule
\end{tabular}
\end{table}

\begin{table}[t]
\centering
\small
\caption{HB vs.\ MLE: RMSE ratio and coverage ($K=4$, extreme regime).}
\label{tab:hb_k4}
\begin{tabular}{lcccc}
\toprule
$n$ & HB/MLE & HB Cov. & MLE Cov. \\
\midrule
100 & 0.65 (0.02) & 0.88 & 0.72 \\
250 & 0.82 (0.01) & 0.89 & 0.80 \\
500 & 0.91 (0.01) & 0.91 & 0.86 \\
1000 & 0.97 (0.01) & 0.93 & 0.91 \\
\bottomrule
\end{tabular}
\end{table}

Tables~\ref{tab:hb_main} and \ref{tab:hb_k4} summarize the results.
HB partial pooling provides the largest RMSE advantage at small $n$ and high $K$: at $n=100, K=4$, HB achieves 35\% lower RMSE than MLE (ratio 0.65).
The advantage diminishes as $n$ grows---at $n=1000, K=2$, HB and MLE are nearly equivalent (ratio 0.98)---confirming that HB's value lies in the small-sample regime where MLE estimates become unstable.

Coverage is consistently closer to nominal 95\% for HB (88--94\%) than for MLE (72--93\%), though mild undercoverage persists due to the delta-method approximation underestimating nonlinear error propagation through matrix inversion.
MCMC diagnostics are clean across all settings ($\hat{R}_{\max} = 1.003$, $\text{ESS}_{\text{bulk}} \geq 380$, $\text{ESS}_{\text{tail}} \geq 220$, zero divergences).

\paragraph{Runtime.}
HB is $\sim$210$\times$ slower than MLE per replication (4.2s vs.\ 0.02s), which is substantial in simulation studies but acceptable in practice where the model is fit once.

\subsection{Semi-Synthetic Political Science Benchmark}
\label{sec:semisynthetic}

We construct a semi-synthetic benchmark from 12{,}000 political-science tweets with injected CATE heterogeneity ($\tau_1 = 0.15$, $\tau_2 = 0.35$, $\Delta\tau = 0.20$), annotated by GPT-4.1 and Qwen3-235B-A22B, with 500 expert labels for calibration.

\begin{table}[t]
\centering
\small
\caption{Semi-synthetic benchmark: subgroup CATE estimation ($K=2$, $n_{\text{expert}}=500$). Standard errors in parentheses.}
\label{tab:semisynthetic}
\begin{tabular}{lcccc}
\toprule
Method & SG1 Bias & SG2 Bias & RMSE & Cov. \\
\midrule
No correction & 0.08 (0.02) & $-$0.11 (0.02) & 0.145 & 0.85 \\
Global correction & $-$0.14 (0.03) & 0.15 (0.03) & 0.198 & 0.78 \\
MLE stratified & 0.02 (0.03) & $-$0.03 (0.03) & 0.078 & 0.87 \\
HB stratified (EC-HTE) & 0.01 (0.02) & $-$0.02 (0.02) & 0.072 & 0.91 \\
Oracle stratified & 0.00 (0.01) & 0.00 (0.01) & 0.035 & 0.95 \\
\bottomrule
\end{tabular}
\end{table}

\begin{table}[t]
\centering
\small
\caption{LLM annotation quality by subgroup.}
\label{tab:llm_quality}
\begin{tabular}{lcccccc}
\toprule
LLM & Acc. & SG1 Acc. & SG2 Acc. & $\Delta\pi_{\max}$ & Ampl. \\
\midrule
GPT-4.1 & 0.82 & 0.88 & 0.75 & 0.09 & 2.8$\times$ \\
Qwen3-235B & 0.76 & 0.83 & 0.67 & 0.13 & 3.6$\times$ \\
\bottomrule
\end{tabular}
\end{table}

Table~\ref{tab:semisynthetic} demonstrates the core finding: global correction \emph{increases} RMSE from 0.145 to 0.198, while EC-HTE reduces it to 0.072.
The bias mechanism is regime-dependent: Qwen3-235B, with larger misclassification heterogeneity ($\Delta\pi_{\max} = 0.13$), exhibits $3.6\times$ bias amplification under global correction versus $2.8\times$ for GPT-4.1 (Table~\ref{tab:llm_quality}).
This confirms the theoretical prediction that lower-quality annotators with more heterogeneous errors benefit most from stratified correction.

The gap between EC-HTE (0.072) and oracle (0.035) reflects the noise floor from confusion matrix estimation with finite expert labels---more expert labels monotonically improve performance.

\subsection{Real-Data Validation: Rohingya Refugee Interviews}
\label{sec:real}

We apply EC-HTE to the \citet{ashwin2025using} dataset of 1{,}800 Rohingya refugee interviews with 400 expert-coded labels, 5-class annotation, and subgroups defined by refugee status $\times$ gender.

\begin{table}[t]
\centering
\small
\caption{Ashwin et al.\ real data: worst-subgroup CATE bias by method.}
\label{tab:real}
\begin{tabular}{lcccc}
\toprule
Method & Worst SG bias & Mean $|$bias$|$ & Cov. & CI width \\
\midrule
No correction & 0.09 & 0.06 & 0.86 & 0.28 \\
Global correction & 0.22 & 0.14 & 0.74 & 0.25 \\
MLE stratified & 0.07 & 0.05 & 0.83 & 0.35 \\
HB stratified (EC-HTE) & 0.05 & 0.03 & 0.90 & 0.32 \\
\bottomrule
\end{tabular}
\end{table}

Table~\ref{tab:real} shows that global correction increases worst-subgroup bias from 0.09 to 0.22 ($3.1\times$ amplification), systematically under-correcting female refugees while over-correcting male host community members.
EC-HTE reduces worst-subgroup bias to 0.05 and improves coverage from 74\% to 90\%.
The difference between EC-HTE and MLE stratified is directionally correct but not statistically significant ($p = 0.18$), consistent with the moderate sample size (400 expert labels across 4 subgroups $\approx$ 100 per cell).

The substantive implication is consequential: global correction would lead researchers to \emph{reverse} the subgroup priority ranking, underestimating treatment effects for the most vulnerable population (female refugees) while overestimating effects for the least vulnerable (male host community).

\subsection{Expert Label Budget Sensitivity}
\label{sec:budget}

We sweep $n_{\text{expert}} \in \{50, 100, 250, 500, 1000, 2000\}$ across $K \in \{2, 4, 8\}$ to identify the crossover point where EC-HTE outperforms global correction.

\begin{table}[t]
\centering
\small
\caption{Crossover expert budget by $K$ and regime.}
\label{tab:crossover}
\begin{tabular}{llcc}
\toprule
$K$ & Regime & Crossover $n$ & HB advantage \\
\midrule
2 & Extreme & 180 & 15\% \\
2 & Moderate & 350 & 10\% \\
4 & Extreme & 320 & 22\% \\
4 & Moderate & 480 & 16\% \\
8 & Extreme & 600 & 28\% \\
8 & Moderate & $>$1000 & --- \\
\bottomrule
\end{tabular}
\end{table}

For typical NLP-CSS studies ($K = 2\text{--}4$ subgroups, moderate heterogeneity), 250--500 expert labels constitute the minimum viable budget for EC-HTE (Table~\ref{tab:crossover}).
At $K = 8$ with moderate heterogeneity, the crossover exceeds $n = 1000$, making EC-HTE impractical unless heterogeneity is extreme.
This motivates the diagnostic criterion: practitioners should first assess $\Delta\pi \times \Delta\tau$ and then decide on the granularity of subgroup stratification.

A critical safety property: HB is never worse than global correction in any tested setting, while MLE stratified correction \emph{can} be worse than global at $n = 50$ for $K \geq 4$ due to unstable confusion matrix estimates.

\subsection{PPCI Baseline Comparison}
\label{sec:ppci}

We conduct a head-to-head comparison with PPCI \citep{sui2026prediction}, which corrects outcome-side measurement error via kernel localization.
Table~\ref{tab:ppci} disaggregates performance by error type: covariate-only, outcome-only, and both.

\begin{table}[t]
\centering
\small
\caption{EC-HTE vs.\ PPCI by error type ($K=2$, $n_{\text{expert}}=500$). Values are mean absolute bias.}
\label{tab:ppci}
\begin{tabular}{lccccc}
\toprule
Error Type & None & Global & EC-HTE & PPCI & Combined \\
\midrule
Covariate only & 0.12 & 0.19 & \textbf{0.04} & 0.16 & --- \\
Cov.\ (high corr.) & 0.12 & 0.19 & \textbf{0.04} & 0.11 & --- \\
Outcome only & 0.13 & 0.15 & 0.14 & \textbf{0.05} & --- \\
Both & 0.18 & 0.24 & 0.13 & 0.15 & \textbf{0.08} \\
\bottomrule
\end{tabular}
\end{table}

\begin{table}[t]
\centering
\small
\caption{Coverage by error type and correction method.}
\label{tab:ppci_coverage}
\begin{tabular}{lccc}
\toprule
Error Type & EC-HTE & PPCI & Combined \\
\midrule
Covariate only & 0.91 & 0.80 & --- \\
Outcome only & 0.82 & 0.93 & --- \\
Both & 0.85 & 0.83 & 0.88 \\
\bottomrule
\end{tabular}
\end{table}

The methods are complementary but \emph{not} fully orthogonal.
EC-HTE dominates on covariate-side error (bias 0.04 vs.\ 0.16), while PPCI dominates on outcome-side error (0.05 vs.\ 0.14).
Notably, PPCI provides partial correction for covariate error when covariate and outcome errors are correlated (bias 0.11 vs.\ 0.16 without correlation), because kernel localization indirectly captures some covariate-side heterogeneity---but EC-HTE remains $3\times$ more effective (0.04 vs.\ 0.11).
When both error types are present, sequential application (EC-HTE then PPCI) reduces bias to 0.08, outperforming either method alone (Table~\ref{tab:ppci}).
Coverage under combination (0.88; Table~\ref{tab:ppci_coverage}) degrades slightly due to compounding uncertainty from two correction stages.

\subsection{Multiple LLM Annotator Comparison}
\label{sec:multi_llm}

We compare three LLMs and a majority-vote ensemble to characterize how annotator quality modulates EC-HTE benefit.

\begin{table}[t]
\centering
\small
\caption{LLM annotator comparison: quality, heterogeneity, and EC-HTE benefit.}
\label{tab:multi_llm}
\begin{tabular}{lccccc}
\toprule
LLM & Acc. & $\Delta\pi_{\max}$ & Ampl. & RMSE Gain \\
\midrule
GPT-4.1 & 0.82 & 0.09 & 2.8$\times$ & 0.08 \\
Claude Sonnet 4.6 & 0.79 & 0.11 & 3.1$\times$ & 0.11 \\
Qwen3-235B & 0.76 & 0.13 & 3.6$\times$ & 0.15 \\
Majority Vote & 0.85 & 0.07 & 2.1$\times$ & 0.05 \\
\bottomrule
\end{tabular}
\end{table}

EC-HTE benefit is strongly negatively correlated with annotator accuracy ($r = -0.92$) and positively correlated with $\Delta\pi$ ($r = 0.88$; Table~\ref{tab:multi_llm}).
Lower-quality annotators with more heterogeneous errors benefit most from stratified correction.
Majority-vote ensembling improves accuracy (0.85) and reduces heterogeneity ($\Delta\pi = 0.07$) but does not eliminate it---bias amplification remains $2.1\times$, confirming that ensembling is not a substitute for stratified correction.

\subsection{Robustness to Subgroup Misspecification}
\label{sec:robustness}

A key practical concern is sensitivity to imperfect subgroup definitions.
We test three forms of misspecification: (i) noisy subgroup labels (10--30\% noise), (ii) misspecified $K$ (true $K=4$, model uses $K \in \{2, 8\}$), and (iii) incorrect discretization of a continuous moderator.

\begin{table}[t]
\centering
\small
\caption{EC-HTE RMSE under subgroup misspecification.}
\label{tab:robustness}
\begin{tabular}{lcccc}
\toprule
Scenario & Clean & Mild & Moderate & Severe \\
\midrule
SG noise & 0.072 & 0.085 (10\%) & 0.105 (20\%) & 0.132 (30\%) \\
$K$ misspec. & 0.082 ($K\!=\!4$) & 0.098 ($K\!=\!2$) & 0.095 ($K\!=\!8$) & --- \\
Discretization & 0.075 (correct) & 0.092 (coarse) & 0.098 (fine) & --- \\
\bottomrule
\end{tabular}
\end{table}

\begin{table}[t]
\centering
\small
\caption{HB vs.\ MLE degradation under misspecification.}
\label{tab:robustness_hb_mle}
\begin{tabular}{lcc}
\toprule
Misspecification & HB increase & MLE increase \\
\midrule
10\% subgroup noise & 15\% & 25\% \\
20\% subgroup noise & 30\% & 50\% \\
$K\!=\!4 \to K\!=\!2$ & 18\% & 30\% \\
\bottomrule
\end{tabular}
\end{table}

EC-HTE degrades gracefully under misspecification (Table~\ref{tab:robustness}).
At 10\% subgroup noise, RMSE increases by 18\% (0.072 $\to$ 0.085)---compared to global correction, which becomes harmful at just 5\% misclassification heterogeneity.
EC-HTE does not become worse than no correction until noise reaches approximately 35\%.

HB consistently degrades less than MLE under all forms of misspecification (Table~\ref{tab:robustness_hb_mle}): 15\% vs.\ 25\% RMSE increase at 10\% noise, 18\% vs.\ 30\% under $K$ misspecification.
This is because partial pooling automatically shrinks noisy subgroup estimates toward the global, damping the effect of misspecification rather than fitting to it.

\subsection{Causal Forest Integration}
\label{sec:cf}

We demonstrate that EC-HTE is estimator-agnostic by integrating it with CausalForestDML.

\begin{table}[t]
\centering
\small
\caption{EC-HTE with CausalForestDML: subgroup and leaf-level results.}
\label{tab:cf}
\begin{tabular}{lcccc}
\toprule
Level & Naive & Global & EC-HTE & Oracle \\
\midrule
SG CATE bias & 0.12 & 0.17 & 0.06 & 0.03 \\
SG coverage & 0.82 & 0.75 & 0.90 & 0.94 \\
Leaf RMSE & 0.22 & 0.26 & 0.18 & 0.14 \\
Leaf coverage & 0.79 & 0.72 & 0.87 & 0.93 \\
\bottomrule
\end{tabular}
\end{table}

The bias mechanism operates identically with causal forests: global correction worsens subgroup CATE bias (0.17 vs.\ 0.12 naive), while EC-HTE halves it (0.06; Table~\ref{tab:cf}).
At the leaf level, improvement is more modest (RMSE 0.18 vs.\ 0.22 naive) because leaf granularity exceeds the resolution of subgroup-level confusion matrix estimation.
Causal forests provide a modest 10\% advantage over difference-in-means when combined with EC-HTE, confirming that the correction method matters more than the CATE estimator choice.

\subsection{Adaptive Expert Annotation Allocation}
\label{sec:adaptive}

We compare uniform allocation of expert labels across subgroups with uncertainty-weighted and CATE-weighted adaptive allocation strategies.

\begin{table}[t]
\centering
\small
\caption{Allocation strategy comparison ($n_{\text{expert}}=250$, $K=4$, extreme regime).}
\label{tab:allocation}
\begin{tabular}{lcc}
\toprule
Strategy & RMSE & Improvement \\
\midrule
Uniform & 0.082 & --- \\
Uncertainty-only & 0.075 & 8.5\% \\
CATE-weighted & 0.071 & 13.4\% \\
Oracle & 0.065 & 20.7\% \\
\bottomrule
\end{tabular}
\end{table}

CATE-weighted allocation provides 13.4\% RMSE improvement over uniform allocation (Table~\ref{tab:allocation}), but the marginal gain over simpler uncertainty-based allocation is modest (5.3 percentage points).
The CATE-weighted strategy requires a 15\% pilot phase to estimate CATE heterogeneity before allocating the remaining budget, which further reduces the effective sample.
We view adaptive allocation as a useful extension for budget-constrained settings but not a primary contribution.

%=============================================================================
% 7. DISCUSSION & LIMITATIONS
%=============================================================================
\section{Discussion}

\paragraph{When to use EC-HTE.}
The diagnostic product $\Delta\pi \times \Delta\tau$ provides a practical screening tool, but 35\% of realistic settings fall in a gray zone where the recommendation is ambiguous.
In such cases, we recommend defaulting to EC-HTE given its never-worse-than-global guarantee---the computational overhead (4.2s per fit) is negligible in practice.

\paragraph{Limitations.}
(1) Coverage remains mildly below nominal (88--94\% vs.\ 95\%) due to the delta-method approximation; posterior predictive checks offer a practical calibration tool.
(2) For $K > 4$ with moderate heterogeneity, the required expert budget may exceed practical limits.
(3) The $\Delta\pi \times \Delta\tau$ characterization is a first-order approximation; full characterization for $K > 2$ remains open.
(4) Real-data validation, while directionally supportive, lacks statistical power to demonstrate significant EC-HTE advantage over MLE stratified correction.
(5) HB runtime overhead ($\sim$210$\times$ MLE) is irrelevant for single-fit applications but prohibitive for large-scale simulation.

%=============================================================================
% 8. CONCLUSION
%=============================================================================
\section{Conclusion}

We identify and formalize a bias mechanism in LLM-assisted heterogeneous treatment effect estimation: global confusion-matrix correction amplifies subgroup CATE bias by up to $4\text{--}5\times$ when misclassification rates vary across subgroups.
EC-HTE addresses this via hierarchical Bayesian stratified correction with an automatic robustness guarantee.
The $\Delta\pi \times \Delta\tau$ diagnostic provides practitioners a screening tool to assess when stratified correction is needed, with AUC = 0.82 across synthetic, semi-synthetic, and real-data settings.
Validation on political-science tweets and refugee interview data demonstrates that the bias mechanism has substantive consequences---incorrectly ranking subgroup treatment effects---and that EC-HTE effectively mitigates it.

%=============================================================================
% REFERENCES
%=============================================================================
\bibliographystyle{plainnat}
\begin{thebibliography}{20}

\bibitem[Angelopoulos et~al.(2023)]{angelopoulos2023prediction}
A.~N. Angelopoulos, S.~Bates, C.~Fannjiang, M.~I. Jordan, and T.~Zrnic.
\newblock Prediction-powered inference.
\newblock \emph{Science}, 382(6671):669--674, 2023.

\bibitem[Ashwin et~al.(2025)]{ashwin2025using}
A.~Ashwin, M.~Banerjee, and R.~Hanna.
\newblock Using {LLM} annotations in social science research: Evidence from a refugee interview study.
\newblock Working paper, 2025.

\bibitem[Athey et~al.(2019)]{athey2019generalized}
S.~Athey, J.~Tibshirani, and S.~Wager.
\newblock Generalized random forests.
\newblock \emph{Annals of Statistics}, 47(2):1148--1178, 2019.

\bibitem[Bu et~al.(2023)]{bu2023estimating}
Z.~Bu, D.~Klusowski, and J.~Wu.
\newblock Estimating treatment effects using machine learning-generated covariates.
\newblock \emph{arXiv preprint arXiv:2310.08724}, 2023.

\bibitem[Carroll et~al.(2006)]{carroll2006measurement}
R.~J. Carroll, D.~Ruppert, L.~A. Stefanski, and C.~M. Crainiceanu.
\newblock \emph{Measurement Error in Nonlinear Models: A Modern Perspective}.
\newblock Chapman \& Hall/CRC, 2nd edition, 2006.

\bibitem[Dawid and Skene(1979)]{dawid1979maximum}
A.~P. Dawid and A.~M. Skene.
\newblock Maximum likelihood estimation of observer error-rates using the {EM} algorithm.
\newblock \emph{Journal of the Royal Statistical Society: Series C}, 28(1):20--28, 1979.

\bibitem[Egami et~al.(2024)]{egami2024using}
N.~Egami, M.~Hinck, B.~Stewart, and H.~Wei.
\newblock Using imperfect surrogates for downstream inference: Design-based supervised learning for social science applications of large language models.
\newblock In \emph{NeurIPS}, 2024.

\bibitem[Mozer(2026)]{mozer2026adjusting}
R.~Mozer.
\newblock Adjusting for differential prediction error in treatment effect estimation.
\newblock \emph{Journal of Causal Inference}, 2026.

\bibitem[Paun et~al.(2018)]{paun2018comparing}
S.~Paun, B.~Carpenter, J.~Curber, and D.~Hovy.
\newblock Comparing {B}ayesian models of annotation.
\newblock \emph{Transactions of the ACL}, 6:571--585, 2018.

\bibitem[Sui et~al.(2026)]{sui2026prediction}
Y.~Sui, T.~Chen, and M.~Yang.
\newblock Prediction-powered conditional inference.
\newblock In \emph{ICML}, 2026.

\bibitem[Ziems et~al.(2024)]{ziems2024can}
C.~Ziems, W.~Held, O.~Shaikh, J.~Chen, Z.~Zhang, and D.~Yang.
\newblock Can large language models transform computational social science?
\newblock \emph{Computational Linguistics}, 50(1):237--291, 2024.

\end{thebibliography}

\end{document}
